% ══════════════════════════════════════════════════════════════════
\chapter{Conclusions and Open Questions}
\label{ch:conclusions}
% ══════════════════════════════════════════════════════════════════

% ──────────────────────────────────────────────────────────────────
\section{Summary of Proven Results}

\begin{enumerate}
  \item \textbf{The gradient trap is structural}
        (Chapter~\ref{ch:gradient_trap}).
        Row centering creates a forward gain of
        $\sqrt{(\pi-1)/\pi} \approx 0.826$ per layer, with a locked
        forward/backward ratio of $0.827$.  This was confirmed
        empirically (single-sample experiments reproduce the same
        pattern; variance sweep shows constant ratio across all
        variance factors).  No scalar variance adjustment can make
        both gains equal to~1 simultaneously.

  \item \textbf{Row centering destroys class structure at depth}
        (Chapter~\ref{ch:geometry}).
        Using $k$-NN accuracy as the geometry metric,
        \texttt{row\_centered\_he\_var\_adj} drops to chance level
        (0.110) by 20 layers.  Standard He stays at 0.639.
        The PCA scatter plots that previously suggested geometry
        preservation were measuring \emph{spread}, not
        \emph{structure}.

  \item \textbf{Forward decay kills both geometry and gradients}
        (Chapters~\ref{ch:gradient_trap}--\ref{ch:geometry}).
        These are the same phenomenon: the $0.826^L$ signal decay
        randomizes activations at depth (destroying class
        information) and creates non-uniform gradients (preventing
        learning).

  \item \textbf{Partial centering and layer-balanced are partial fixes}
        (Chapter~\ref{ch:fixing_trap}).
        The alpha sweep reveals a Pareto frontier with no escape.
        The layer-balanced approach ($\eta = 0.5$) reduces gradient
        ratio from $11\times$ to $4\times$ while keeping full geometry,
        but the backward gain of 1.21 limits further improvement.

  \item \textbf{Negative bias cannot fix the kernel}
        (Section~\ref{sec:biased_relu}).
        Optimal $\beta^* \approx -0.93$ reduces distortion by only 13\%,
        at the cost of 82\% silent neurons and $5.7\times$ variance
        compensation.  No pointwise nonlinearity can fully preserve
        the input kernel.

  \item \textbf{The kernel-preserving optimizer confirms centering is
        wrong} (Section~\ref{sec:kp_analysis}).
        The optimizer discovers the \emph{opposite} of row centering
        (large row sums, inflated weights, anisotropic structure), but
        this solution is non-composable across layers.

  \item \textbf{Two different collapse modes}
        (Chapter~\ref{ch:angle_map}).
        He collapses to $0^\circ$ (identity: all vectors align).
        Row-centered collapses to $71^\circ$ (equidistance: all
        vectors become equidistant).
        Both destroy class information, but He preserves
        \emph{ordering} (which $k$-NN needs) while row-centered
        destroys it.

  \item \textbf{Orthogonal\_he is the practical winner}
        (Section~\ref{sec:orthogonal_analysis}).
        It achieves $k$-NN $= 0.662$ at 20 layers (vs.\ He's 0.639),
        a $\sim\!3.6\%$ improvement from reduced variance, with no
        gradient trap.
\end{enumerate}

% ──────────────────────────────────────────────────────────────────
\section{The Real Enemy: $D(\alpha) > 0$}

The systematic ReLU kernel distortion
\[
  D(\alpha) = \frac{1}{\pi}\bigl(\sin\alpha - \alpha\cos\alpha\bigr) > 0
  \quad\text{for all } \alpha \in (0,\pi)
\]
affects \textbf{all} initializers.  No single-layer weight structure
can eliminate it:
\begin{itemize}
  \item The distortion depends on the normalized kernel
        $\hat{K}(\alpha) = K(\alpha)/K(0)$, which is invariant to
        weight scaling.
  \item Orthogonal matrices have the same marginal distribution as
        Gaussian, so the same $\hat{K}$.
  \item Row centering changes the effective kernel but introduces worse
        problems (forward decay, equidistance collapse).
  \item Negative bias reshapes $D_\beta(\alpha)$ but cannot make it
        zero everywhere.
  \item The kernel-preserving optimizer finds a complex solution that
        works per-layer but is non-composable.
\end{itemize}

This distortion is \textbf{fundamental to any pointwise nonlinearity}:
the operation
$\E[\phi(\bm{w}^\top \bm{x}_i) \cdot \phi(\bm{w}^\top \bm{x}_j)]$
cannot equal $\langle \bm{x}_i, \bm{x}_j \rangle$ for all angles
unless $\phi$ is linear.

% ──────────────────────────────────────────────────────────────────
\section{Open Questions}

\begin{enumerate}
  \item \textbf{Reframing row centering in the thesis.}
        The advisor proposed row centering specifically for geometry
        preservation.  Our results show it is worse than standard He
        for class separability.  How should we present this finding
        diplomatically while being scientifically rigorous?

  \item \textbf{Pivoting toward bounding the distortion.}
        If the kernel distortion cannot be eliminated, perhaps the
        thesis contribution is proving this impossibility and
        characterizing the best achievable bounds.  What is the
        optimal single-layer angle map achievable by any weight
        structure?

  \item \textbf{Multi-layer or joint optimization.}
        The kernel-preserving optimizer fails because each layer is
        optimized independently.  Could jointly optimizing all layers
        (accounting for cumulative kernel drift) produce a composable
        solution?

  \item \textbf{Alternative activation functions.}
        The distortion $D(\alpha) > 0$ is specific to ReLU.
        Activation functions with $\E[\phi(z)] = 0$ (e.g., centered
        ReLU, GELU minus its mean, or antisymmetric activations) would
        eliminate the DC problem that drives the row-centering failure.
        Do they also reduce the kernel distortion?

  \item \textbf{Practical implications.}
        For networks trained with gradient descent (not just random
        projections), does the initialization-time geometry matter?
        The gradient trap prevents early learning, but adaptive
        optimizers (Adam, LARS) may compensate.
\end{enumerate}
